\section{Properties of Symmetric Matrices}

% In each part \(P\) and \(Q\) are symmetric.

\textbf{(a)}
\textbf{True}. If \(P\ge 0\), then
\[
(P+Q)-Q=P\ge 0
\]
Therefore
\[
P+Q\ge Q
\]

\textbf{(b)}
\textbf{True}. If \(P\ge Q\), then simply multiply both sides by $-1$
\[
-P\le -Q
\]

% \[
% P-Q\ge 0
% \]
% But
% \[
% (-Q)-(-P)=P-Q
% \]
% so
% \[
% -Q\ge -P
% \]
% which is the same as
% \[
% -P\le -Q
% \]

\textbf{(c)}
\textbf{True}. Let \(P>0\), and let \(x\ne 0\). Since \(P\) is invertible, \(y=P^{-1}x\) is also nonzero. Then
\[
x^T P^{-1}x
=
(Py)^T P^{-1}(Py)
=
y^T P y
>0
\]
so
\[
P^{-1}>0
\]


%  \[
%   P=Q\Lambda Q^T,\qquad \lambda_i(P)>0
%   \]

%   Therefore,

%   \[
%   P^{-1}=Q\Lambda^{-1}Q^T
%   \]

%   so the eigenvalues of $(P^{-1})$ are

%   \[
%   \lambda_i(P^{-1})=\frac{1}{\lambda_i(P)}>0
%   \]

%   Hence (P^{-1}>0).

\textbf{(d)}
\textbf{True}. Assume
\[
P\ge Q>0
\]
Define
\[
Z=Q^{-1/2}P Q^{-1/2}
\]
Then
\[
Z-I=Q^{-1/2}(P-Q)Q^{-1/2}\ge 0
\]
so
\[
Z\ge I
\]
The eigenvalues of \(Z\) are at least one, so the eigenvalues of \(Z^{-1}\) are at most one. Therefore
\[
Z^{-1}\le I
\]
Now
\[
P^{-1}=Q^{-1/2}Z^{-1}Q^{-1/2},
\qquad
Q^{-1}=Q^{-1/2}IQ^{-1/2}
\]
and hence
\[
Q^{-1}-P^{-1}
=
Q^{-1/2}(I-Z^{-1})Q^{-1/2}\ge 0
\]
so
\[
P^{-1}\le Q^{-1}
\]

\textbf{(e)}
\textbf{False}. Take
\[
P=
\begin{bmatrix}
1&0\\
0&1
\end{bmatrix},
\qquad
Q=
\begin{bmatrix}
-2&0\\
0&0
\end{bmatrix}
\]
Then
\[
P-Q=
\begin{bmatrix}
3&0\\
0&1
\end{bmatrix}\ge 0
\]
so \(P\ge Q\). But
\[
P^2-Q^2
=
\begin{bmatrix}
1&0\\
0&1
\end{bmatrix}
-
\begin{bmatrix}
4&0\\
0&0
\end{bmatrix}
=
\begin{bmatrix}
-3&0\\
0&1
\end{bmatrix}
\]
which is not positive semidefinite. Therefore
\[
P^2\not\ge Q^2
\]
